\chapter{Ciągi spektralne: od filtracji do fibracji Serre'a}
\label{ch:ciagi-spektralne-v5}

\section{Motywacja: dlaczego jeden długi ciąg dokładny nie wystarcza?}
Dla fibracji $F\to E\xrightarrow{p}B$ długi ciąg grup homotopii
porównuje $\pi_*(F)$, $\pi_*(E)$ i $\pi_*(B)$, ale nawet
znajomość wszystkich grup $H_*(F)$ i $H_*(B)$ nie pozwala po prostu
pomnożyć ich i otrzymać $H_*(E)$. W fibracji Hopfa
$S^1\to S^3\to S^2$ naiwne mnożenie homologii dawałoby
homologię $S^1\times S^2$, a nie $S^3$. Musi istnieć
mechanizm, który \emph{usuwa} niektóre klasy przy sklejaniu.

Takim mechanizmem jest ciąg spektralny: na każdej \emph{stronie}
$E^r$ widzimy grupy podobne do kolejnych przybliżeń homologii,
a różniczka $d^r$ mierzy obstrukcję przeniesienia klasy do
następnego przybliżenia. Słowo ,,spektralny'' ma tu tradycyjne
znaczenie i nie oznacza widma operatora.

Najpierw wyprowadzimy konstrukcję z filtracji \emph{dowolnego}
kompleksu łańcuchowego. Dopiero potem zastosujemy ją do
filtracji przestrzeni całkowitej wiązki przez przeciwobrazy
szkieletów podstawy. Wszystkie homologie do końca tego rozdziału
mają współczynniki całkowite, chyba że zaznaczono inaczej.

\section{Filtracja kompleksu i pierwsza strona}
\begin{definicja}[Filtracja rosnąca kompleksu]
Niech $(C_*,\partial)$ będzie kompleksem łańcuchowym grup abelowych
w stopniach nieujemnych. \emph{Filtracją rosnącą} nazywamy
podkompleksy
$\cdots\subseteq F_{p-1}C_*\subseteq F_pC_*\subseteq\cdots$,
takie że $\partial(F_pC_n)\subseteq F_pC_{n-1}$.
Filtracja jest \emph{ograniczona w każdym stopniu}, jeżeli
dla każdego $n$ istnieją $a_n\le b_n$ z
$F_{a_n-1}C_n=0$ oraz $F_{b_n}C_n=C_n$.
Dla $p<0$ w przykładach pierwszoćwiartkowych przyjmujemy $F_pC=0$.
\end{definicja}
\begin{przyklad}[Filtracja przez szkielety]
Jeżeli $X$ jest CW-kompleksem, to jego szkielety
$X^{-1}=\varnothing\subset X^0\subset X^1\subset\cdots$
dają filtrację komórkowego kompleksu $C_*^{\mathrm{CW}}(X)$.
W każdej grupie łańcuchów granica komórki $n$-wymiarowej
leży w szkielecie $n-1$.
\end{przyklad}
\begin{definicja}[Strona zerowa i pierwsza]
Dla filtrowanego kompleksu $C$ definiujemy
\[
E^0_{p,q}=F_pC_{p+q}/F_{p-1}C_{p+q},\qquad
 d^0:E^0_{p,q}\to E^0_{p,q-1},\quad [c]\mapsto[\partial c].
\]
Jest to dobrze określone, a $(d^0)^2=0$.
\emph{Pierwsza strona} to $E^1_{p,q}=H_{p+q}(F_pC/F_{p-1}C)$.
Grupa $E^1_{p,q}$ jest grupą homologii ilorazu kompleksów,
a nie ilorazem $H(F_pC)/H(F_{p-1}C)$.
\end{definicja}
\begin{przyklad}[Elementarna filtracja]
Niech $C_2=\Z x$, $C_1=\Z y$, $\partial x=y$,
pozostałe grupy są zerowe, oraz $F_0C=\Z y$,
$F_1C=F_0C$, $F_2C=C$. Zwykła homologia $C$ znika,
ale pierwsza strona ma $E^1_{2,0}=\Z$ i $E^1_{0,1}=\Z$.
Różniczka $d^1$ jest zerowa ze względów stopniowania,
więc oba elementy przeżywają do strony drugiej.
Dopiero $d^2:E^2_{2,0}\to E^2_{0,1}$ jest izomorfizmem.
To najprostszy przykład klasy, która ,,ginie dopiero na drugiej stronie''.
\end{przyklad}

\begin{przyklad}[Różniczka pojawiająca się dopiero na trzeciej stronie]
\label{prz:spectral-d3-cancel}
Ustalmy kompleks $C_3=\Z x$, $C_2=\Z y$,
$\partial_3x=y$, z pozostałymi grupami zerowymi.
Jego homologia jest zerowa.
Zdefiniujmy filtrację
\[
 F_{-1}C=0,\qquad F_0C=\Z y,\qquad
 F_1C=F_2C=F_0C,\qquad F_3C=C.
\]
Na stronie pierwszej występują jedynie
$E^1_{3,0}\cong\Z$ (klasa $x$)
i $E^1_{0,2}\cong\Z$ (klasa $y$).
Obie różniczki $d^1$ i $d^2$ muszą znikać:
dla $x$ cel miałby pierwszy indeks $2$ albo $1$,
gdzie nie ma grup niezerowych;
dla $y$ nie ma celu w pierwszej ćwiartce.
Natomiast
\[
 d^3:E^3_{3,0}\longrightarrow E^3_{0,2},
 \qquad [x]\longmapsto[\partial x]=[y]
\]
jest izomorfizmem. Strona czwarta jest zerowa,
zgodnie z $H_*(C)=0$.
Ta sama konstrukcja z $C_r=\Z x$,
$C_{r-1}=\Z y$ i filtrem
$F_0C=\Z y$, $F_{r-1}C=F_0C$, $F_rC=C$
pokazuje dla każdego $r\ge2$ przykład
pierwszej niezerowej różniczki \emph{właśnie} $d^r$.
\end{przyklad}

\section{Konstrukcja wszystkich różniczek i dowód strony następnej}
\begin{definicja}[Cykle i brzegi rzędu $r$]
Niech $r\ge1$ i $n=p+q$. Określamy
\[
Z^r_{p,q}=\{x\in F_pC_n:\partial x\in F_{p-r}C_{n-1}\},
\quad
B^r_{p,q}=Z^r_{p,q}\cap\bigl(F_{p-1}C_n+\partial(F_{p+r-1}C_{n+1})\bigr).
\]
Kładziemy $E^r_{p,q}=Z^r_{p,q}/B^r_{p,q}$.
Dla $r=1$ identyfikujemy ten iloraz z poprzednio zdefiniowaną
pierwszą stroną.
\end{definicja}
\begin{lemat}[Poprawność różniczki]
\label{lem:dr-v5}
Dla $r\ge1$ wzór $d^r[x]=[\partial x]$ wyznacza
homomorfizm
$d^r:E^r_{p,q}\to E^r_{p-r,q+r-1}$,
a $d^r\circ d^r=0$.
\end{lemat}
\begin{proof}
Niech $x\in Z^r_{p,q}$. Wtedy $\partial x\in F_{p-r}C_{n-1}$,
a $\partial^2x=0$, zatem $\partial x$ jest cyklem
rzędu $r$ w miejscu $(p-r,q+r-1)$. Załóżmy, że zmienimy
reprezentanta na $x'=x+y+\partial z$, gdzie
$y\in F_{p-1}C_n$, $z\in F_{p+r-1}C_{n+1}$,
a $x'\in Z^r_{p,q}$. Wówczas
$\partial x'-\partial x=\partial y$,
co jest brzegiem rzędu $r$ w miejscu docelowym, bo
$y\in F_{(p-r)+r-1}C_n=F_{p-1}C_n$.
Własność $(d^r)^2=0$ pochodzi z $\partial^2=0$.
\end{proof}
\begin{twierdzenie}[Strony sukcesywne]
\label{thm:pages-v5}
Dla każdego $r\ge1$ istnieje naturalny izomorfizm
$E^{r+1}_{p,q}\cong H_{p,q}(E^r,d^r)$, czyli
$E^{r+1}_{p,q}\cong\ker(d^r:E^r_{p,q}\to E^r_{p-r,q+r-1})/
\operatorname{im}(d^r:E^r_{p+r,q-r+1}\to E^r_{p,q})$.
\end{twierdzenie}
\begin{proof}
Wyjaśnijmy osobno przechodzenie do jądra i następnie
ilorazowanie przez obraz. Warunek $d^r[x]=0$ mówi, że
$\partial x$ można zapisać jako
$\partial x=y+\partial z$, gdzie
$y\in F_{p-r-1}C_{n-1}$ i $z\in F_{p-1}C_n$.
Zastąpmy $x$ przez $x-z$ (wolno, gdyż $z\in F_{p-1}$).
Wtedy $\partial(x-z)=y\in F_{p-r-1}C_{n-1}$,
więc $x-z\in Z^{r+1}_{p,q}$. Odwrotnie, cykl rzędu $r+1$
jest cyklem rzędu $r$ zabijanym przez $d^r$.

Ustalmy teraz dwie reprezentacje $x,x'\in Z^{r+1}_{p,q}$.
Różnica ich klas na stronie $r$ jest obrazem $d^r$
wtedy i tylko wtedy, gdy modulo $F_{p-1}C_n$ jest
brzegiem elementu z $F_{p+r}C_{n+1}$, którego brzeg
leży w $F_pC_n$; ostatni warunek jest dokładnie
warunkiem bycia cyklem rzędu $r$ w miejscu $(p+r,q-r+1)$.
Otrzymujemy więc relację
$x-x'\in F_{p-1}C_n+\partial(F_{p+r}C_{n+1})$,
z uwzględnieniem warunku bycia w $Z^{r+1}$.
Jest to definicja mianownika $B^{r+1}_{p,q}$.
Opisane przejście jest zgodne z dodawaniem i daje izomorfizm.
\end{proof}
\begin{uwaga}[Indeksy nie są dekoracją]
Różniczka w \emph{homologicznym} ciągu spektralnym ma
bistopień $(-r,r-1)$. Dla \emph{kohomologicznego} ciągu
spektralnego przy standardowej konwencji ma bistopień
$(r,1-r)$. W dalszym rachunku każdorazowo zaznaczamy,
którą konwencję stosujemy.
\end{uwaga}

\section{Zbieżność i problem rozszerzeń}
\begin{definicja}
Filtracja kompleksu indukuje filtrację homologii przez
$F_pH_n(C)=\operatorname{im}(H_n(F_pC)\to H_n(C))$.
Piszemy $\operatorname{gr}_pH_n(C)=F_pH_n(C)/F_{p-1}H_n(C)$.
\end{definicja}
\begin{twierdzenie}[Zbieżność przy filtracji ograniczonej]
\label{thm:convergence-v5}
Niech $F_pC$ będzie filtracją ograniczoną w każdym stopniu.
Wtedy dla każdego $(p,q)$ grupy $E^r_{p,q}$ stabilizują się,
a ich wartość stabilna $E^\infty_{p,q}$ jest naturalnie
izomorficzna z $\operatorname{gr}_pH_{p+q}(C)$.
\end{twierdzenie}
\begin{proof}
Dla ustalonego $n=p+q$, gdy $r$ jest dostatecznie duże,
$F_{p-r}C_{n-1}=0$ oraz $F_{p+r-1}C_{n+1}=C_{n+1}$.
Dlatego licznik w definicji $E^r$ staje się
$Z_n(C)\cap F_pC_n$, zaś mianownik
$\bigl(Z_n(C)\cap F_pC_n\bigr)\cap
(F_{p-1}C_n+\partial C_{n+1})$.
Przyporządkowanie cyklowi $z\in Z_n(C)\cap F_pC_n$
jego klasy $[z]\in H_n(C)$ ma obraz równy $F_pH_n(C)$.
Jego przeciwobraz $F_{p-1}H_n(C)$ składa się dokładnie
z cykli $z=y+\partial w$, gdzie
$y\in F_{p-1}C_n$ i $w\in C_{n+1}$.
Ponieważ $\partial z=0$ oraz $\partial^2w=0$,
również $\partial y=0$; odwrotnie, dowolna klasa z
$F_{p-1}H_n(C)$ ma takiego reprezentanta $y$.
Jądro odwzorowania do ilorazu przez $F_{p-1}H_n(C)$
jest więc dokładnie stabilnym mianownikiem.
Na mocy twierdzenia o izomorfizmie otrzymujemy
$F_pH_n(C)/F_{p-1}H_n(C)$. Ponieważ filtry w danym
stopniu są ograniczone, nie występuje problem nieskończonych
kroków ani dopełniania granic.
\end{proof}
\begin{przyklad}[Nie wolno sumować $E^\infty$ bez uzasadnienia]
Jeżeli w pewnym stopniu filtracja ma dwa kolejne ilorazy
$E^\infty_{0,n}\cong\Z/2$ i
$E^\infty_{1,n-1}\cong\Z/2$, to grupa $H_n$ może być
izomorficzna zarówno z $\Z/4$, jak i z
$\Z/2\oplus\Z/2$: obie mają filtrację z tymi ilorazami.
Dane $E^\infty$ same nie rozwiązują \emph{problemu rozszerzeń}.
Jeśli natomiast wszystkie występujące ilorazy są grupami wolnymi,
to kolejne krótkie ciągi dokładne rozszczepiają się jako ciągi
grup abelowych, choć zwykle niekanonicznie.
\end{przyklad}
\begin{przyklad}[Dwa filtrowane kompleksy, ta sama strona stabilna]
\label{prz:spectral-extension-z4}
Skonstruujemy dokładnie dwie filtracje ilustrujące
poprzedni problem rozszerzeń. Oba kompleksy mają
zerowe różniczki, a jedyną niezerową grupę
łańcuchów w stopniu $1$.
W pierwszym przypadku kładziemy
\[
 C_1=\Z/4\Z,\qquad
 F_{-1}C_1=0,\quad
 F_0C_1=2\Z/4\Z,\quad
 F_1C_1=C_1.
\]
W drugim bierzemy
\[
 C'_1=(\Z/2\Z)\oplus(\Z/2\Z),\qquad
 F_{-1}C'_1=0,\quad
 F_0C'_1=(\Z/2\Z)\oplus0,\quad
 F_1C'_1=C'_1.
\]
Filtracje w innych stopniach są zerowe; dla
$p\ge1$ pozostają równe całym grupom łańcuchów.
W obu przypadkach na stronie pierwszej są tylko
\[
 E^1_{0,1}\cong\Z/2\Z,\qquad
 E^1_{1,0}\cong\Z/2\Z.
\]
Wszystkie różniczki $d^r$ znikają, gdyż pochodzą
od zerowej różniczki kompleksu wyjściowego.
Stąd strony $E^1,E^2,\ldots,E^\infty$
są \emph{identyczne jako grupy bistopniowane}.
W pierwszym przypadku jednak
$H_1(C)=\Z/4\Z$, a w drugim
$H_1(C')=(\Z/2\Z)^2$.
W obu otrzymujemy dokładny ciąg
\[
 0\longrightarrow\Z/2\Z
 \longrightarrow H_1
 \longrightarrow\Z/2\Z
 \longrightarrow0,
\]
ale tylko w drugim rozszczepia się on zgodnie
z podaną filtracją. Ponieważ w $\Z/4\Z$
istnieje element rzędu $4$, a w $(\Z/2\Z)^2$
każdy niezerowy element ma rząd $2$,
grupy homologii nie są izomorficzne.
Nie wystarczy więc poznać nawet \emph{wszystkich}
stron i różniczek — trzeba dodatkowo rozwiązać
problem rozszerzenia filtracji.
\end{przyklad}

\begin{figure}[!htbp]
\centering
\begin{tikzpicture}[>=Stealth,font=\small]
 \node[draw=manifoldblue!75!black,rounded corners=4pt,
       fill=manifoldblue!8,minimum width=4.0cm,
       minimum height=1.25cm,align=center] (a) at (0,0)
       {$E^\infty_{0,1}=\Z/2\Z$\\[3pt]
        $E^\infty_{1,0}=\Z/2\Z$};
 \node[draw=manifoldgreen!70!black,rounded corners=4pt,
       fill=manifoldgreen!8,minimum width=3.1cm,
       minimum height=1.0cm,align=center] (b) at (-2.15,-2.20)
       {$H_1=\Z/4\Z$\\[2pt] nierozszczepiony};
 \node[draw=manifoldgold!80!black,rounded corners=4pt,
       fill=manifoldgold!10,minimum width=3.6cm,
       minimum height=1.0cm,align=center] (c) at (2.15,-2.20)
       {$H_1=(\Z/2\Z)^2$\\[2pt] rozszczepiony};
 \draw[->,thick] (a.south west)--(b.north);
 \draw[->,thick] (a.south east)--(c.north);
\end{tikzpicture}
\caption{Identyczne ilorazy filtracji $E^\infty$
nie określają, czy grupa homologii jest
$\Z/4\Z$, czy $(\Z/2\Z)^2$.}
\label{fig:spectral-extension-z4}
\end{figure}

\begin{uwaga}
Przy filtracjach nieograniczonych trzeba osobno sformułować
warunki zbieżności i kompletności. Nie używamy
Twierdzenia~\ref{thm:convergence-v5} bez tych hipotez.
\end{uwaga}

\section{Dokładna para: skąd pochodzą strony z punktu widzenia par}
\begin{definicja}[Dokładna para]
Dokładną parą grup abelowych nazywamy układ
$D\xrightarrow{i}D\xrightarrow{j}E\xrightarrow{k}D$
z $\operatorname{im}i=\ker j$,
$\operatorname{im}j=\ker k$,
$\operatorname{im}k=\ker i$.
Wtedy $d=j\circ k:E\to E$ spełnia $d^2=0$.
\end{definicja}
\begin{proof}
Ponieważ $kj=0$ z dokładności w środkowym wierzchołku,
$d^2=j(kj)k=0$.
\end{proof}
\begin{stwierdzenie}
Filtracja kompleksu $F_pC$ wyznacza dokładną parę
$D^1_{p,q}=H_{p+q}(F_pC)$ oraz
$E^1_{p,q}=H_{p+q}(F_pC,F_{p-1}C)$.
Przy tych indeksach mapy mają postać
\[
 i:D^1_{p-1,q+1}\to D^1_{p,q},\qquad
 j:D^1_{p,q}\to E^1_{p,q},\qquad
 k:E^1_{p,q}\to D^1_{p-1,q}.
\]
Pochodzą odpowiednio z włączenia filtru,
przejścia do ilorazu i homomorfizmu łączącego
w długim ciągu dokładnym pary kompleksów.
Przy tej konstrukcji $d^1=j\circ k$, a kolejne pary
pochodne dają strony $E^r$ z poprzedniej sekcji.
\end{stwierdzenie}
\begin{proof}
Dla każdego $p$ stosujemy krótki ciąg dokładny kompleksów
$0\to F_{p-1}C\to F_pC\to F_pC/F_{p-1}C\to0$.
Z dokładności jego długiego ciągu wynika kolejno
$\operatorname{im}i=\ker j$, $\operatorname{im}j=\ker k$
i $\operatorname{im}k=\ker i$, gdy uwzględnimy przesunięcie
indeksu $p\mapsto p-1$ mapy $k$.
Na poziomie cyklu $x\in F_pC$ z $\partial x\in F_{p-1}C$
homomorfizm łączący wysyła $[x]$ w $[\partial x]$;
a następnie $j$ bierze jego klasę modulo $F_{p-2}C$.
To jest dokładnie definicja $d^1$.
\end{proof}

\section{Fibracja nad kompleksem CW i lokalne współczynniki}
Długi ciąg homotopii nie wyznacza na ogół homologii $H_*(E)$, w szczególności jej torsji ani iloczynów w kohomologii przestrzeni całkowitej. Nie oznacza to, że grupy homotopii same nie mogą mieć torsji.
Zastosujemy właśnie zbudowany mechanizm do fibracji.
Najpierw podajemy szczegółowy dowód dla \emph{lokalnie trywialnych}
wiązek nad skończeniewymiarowymi CW-kompleksami. Ich rzutowania
są fibracjami Hurewicza w zwykłych sytuacjach CW. Istnieje też
pełne twierdzenie dla fibracji Serre'a; wymaga jednak osobnego
argumentu porównującego homologię względną przeciwobrazów komórek
z homologią ich włókien i nie utożsamiamy go z lokalną trywialnością.

\begin{lemat}[Fibracja nad dyskiem]
\label{lem:serre-disk-v5}
Niech $p:E\to D^m$ będzie fibracją Hurewicza, $b_0\in D^m$
i $F=p^{-1}(b_0)$. Wówczas $p$ jest homotopijnie
trywialna \emph{nad podstawą}: istnieją mapy
$E\rightleftarrows D^m\times F$ nad $D^m$,
których złożenia są homotopijne (nad $D^m$) do identyczności.
\end{lemat}
\begin{proof}
Wybierzmy kontrakcję podstawy $h_t:D^m\to D^m$,
$h_0=\id$, $h_1\equiv b_0$. Podnosimy homotopię
$H(e,t)=h_t(p(e))$ z podniesienia początkowego $\id_E$.
Otrzymujemy $L:E\times I\to E$ z $pL(e,t)=h_t(p(e))$
i $L(e,0)=e$. W szczególności $r(e)=L(e,1)$ leży w $F$.
Mapa $g:E\to D^m\times F$ dana przez $e\mapsto(p(e),r(e))$
jest więc ciągła. Własność podnoszenia dla
$Z=D^m\times F$ zastosowana do homotopii
$K((b,x),t)=h_{1-t}(b)$ i mapy początkowej
$(b,x)\mapsto x\in F\subset E$ dostarcza mapy
$f:D^m\times F\to E$ nad $D^m$ przy $t=1$.
Złączenie podstawowych kontrakcji i ich homotopii
względem końców daje homotopie włókniste
$fg\simeq\id_E$ i $gf\simeq\id_{D^m\times F}$:
wymagana homotopia podniesień otrzymywana jest przez
własność Hurewicza zastosowaną do rodziny parametryzowanej
przez $E\times I$ i odpowiednio $F\times D^m\times I$.
Można również bezpośrednio sprawdzić ten wniosek przez
podnoszenie liniowej homotopii między odpowiednimi
złożeniami dróg w wypukłym $D^m$. Ważne, że homotopie
pozostają nad podstawą, a nie tylko w przestrzeni całkowitej.
\end{proof}

\begin{definicja}[Monodromia na homologii włókna]
Niech $p:E\to B$ będzie fibracją Hurewicza o łukowo spójnej
podstawie, $b_0\in B$, $F=p^{-1}(b_0)$.
Transport wzdłuż drogi $\gamma:b_0\to b_1$ daje określoną
z dokładnością do homotopii równoważność
$T_\gamma:F_{b_0}\to F_{b_1}$.
W każdym stopniu $q$ otrzymujemy system lokalny
$\mathcal H_q(F)$ o włóknach $H_q(F_b;\Z)$ i transporcie
$(T_\gamma)_*$. Przy zamkniętej pętli działa
$\pi_1(B,b_0)$ na $H_q(F;\Z)$.
\end{definicja}

\begin{lemat}[Względna homologia przeciwobrazów szkieletów]
\label{lem:serre-cells-v5}
Niech $p:E\to B$ będzie lokalnie trywialną wiązką nad
CW-kompleksem $B$, a $E_j=p^{-1}(B^j)$.
Dla $j\ge0$ i $k\in\Z$ istnieje izomorfizm
\[
H_k(E_j,E_{j-1};\Z)\cong
\bigoplus_{\alpha\in\operatorname{Cells}_j(B)}H_{k-j}(F_\alpha;\Z),
\]
gdzie $F_\alpha$ jest włóknem nad wnętrzem komórki
$e^j_\alpha$, a $H_l(F_\alpha)=0$ dla $l<0$.
Dla komórek zerowymiarowych stosujemy konwencję
$D^0=\{*\}$ oraz $S^{-1}=\varnothing$.
Znak identyfikacji zależy od wyboru orientacji komórki.
\end{lemat}
\begin{proof}
CW-kompleks $B^j$ jest otrzymany z $B^{j-1}$ przez
dołączenie dysków charakterystycznych
$\chi_\alpha:D^j_\alpha\to B^j$. Ich otwarte wnętrza
są parami rozłączne. Cofnięcie lokalnie trywialnej wiązki
na dysk $D^j_\alpha$ jest trywialne (wiązka nad
parazwartą przestrzenią ściągalną jest trywialna), więc
nad nim mamy parę
$(D^j\times F_\alpha,S^{j-1}\times F_\alpha)$.
Jej brzeg ma kołnierz: para jest kofibracyjna.

Relatywne wycinanie dla komórkowego przylegania daje
izomorfizm z sumą bezpośrednią względnych homologii
kolejnych par dyskowych. Dla skończonej rodziny komórek
jest to zwykłe wycinanie po pogrubieniu brzegu kołnierzem;
w przypadku nieskończenie wielu komórek stosujemy fakt,
że każdy singularny łańcuch jest skończony, a jego obraz
zwarty, więc obliczenie przechodzi do granicy prostej
po skończonych podkompleksach. Stąd otrzymujemy sumę
prostą (a nie iloczyn).

Grupy $H_l(D^j,S^{j-1};\Z)$ wynoszą $\Z$ dla $l=j$
i zero dla pozostałych $l$. Z relatywnego wzoru
Künnetha wynika zatem
$H_k(D^j\times F,S^{j-1}\times F;\Z)\cong H_{k-j}(F)$;
człon $\operatorname{Tor}$ znika z powodu wolności
homologii pary $(D^j,S^{j-1})$.
Sumując po komórkach, otrzymujemy tezę.
\end{proof}

\begin{twierdzenie}[Homologiczny ciąg spektralny Serre'a]
\label{thm:serre-v5}
Niech $p:E\to B$ będzie lokalnie trywialną wiązką o typie
CW-kompleksu jako włóknie, $B$ spójnym i skończeniewymiarowym
CW-kompleksem, $b_0\in B$, $F=p^{-1}(b_0)$.
Istnieje naturalny pierwszoćwiartkowy ciąg spektralny
\[
E^2_{p,q}\cong H_p(B;\mathcal H_q(F)),\qquad
 d^r:E^r_{p,q}\longrightarrow E^r_{p-r,q+r-1},
\]
zbieżny do $H_{p+q}(E;\Z)$ w sensie skończonej w każdym
stopniu filtracji $F_pH_{p+q}(E)$ z
$E^\infty_{p,q}\cong F_pH_{p+q}(E)/F_{p-1}H_{p+q}(E)$.
Jeżeli monodromia na $H_q(F;\Z)$ jest trywialna, można
zastąpić $\mathcal H_q(F)$ stałym systemem współczynników.
Dla fibracji Serre'a o odpowiednim typie CW analogiczna
konkluzja pozostaje prawdziwa; nie traktujemy jej jednak
jako udowodnionej w pełnej ogólności poniższym argumentem.
\end{twierdzenie}
\begin{proof}
\emph{Krok 1: filtrujemy przestrzeń całkowitą.}
Niech $B^p$ oznacza szkielet $p$-wymiarowy $B$, przyjmijmy
$B^{-1}=\varnothing$ i $E_p=p^{-1}(B^p)$.
Filtrujemy singularny kompleks $C_*(E)$ podkompleksami
$F_pC_*(E)=\operatorname{im}(C_*(E_p)\to C_*(E))$.
Każda singularna mapa sympleksu do $E$ ma zwarty obraz,
a jej rzut do CW-kompleksu jest zawarty w skończonym
podkompleksie $B$. Jeżeli $m=\dim B<\infty$, to
$E_m=E$, więc filtracja $C_*(E)$ jest rzeczywiście
ograniczona w każdym stopniu: $F_{-1}C_*=0$ i
$F_mC_*=C_*$. Dzięki temu wolno zastosować
Twierdzenie~\ref{thm:convergence-v5} bez ukrytych
założeń o zupełności filtracji.

\emph{Krok 2: obliczamy wyrazy strony pierwszej.}
Mamy $E^1_{p,q}=H_{p+q}(E_p,E_{p-1};\Z)$.
Wybierzmy dla każdej zorientowanej $p$-komórki $e^p_\alpha$
mapę charakterystyczną
$\chi_\alpha:(D^p,S^{p-1})\to(B^p,B^{p-1})$.
Lemat~\ref{lem:serre-cells-v5} i lokalne trywializacje nad komórkami dają
\[
H_{p+q}(E_p,E_{p-1};\Z)\cong
\bigoplus_{\alpha\in\operatorname{Cells}_p(B)}
H_{p+q}(D^p\times F,S^{p-1}\times F;\Z).
\]
Ponieważ $H_k(D^p,S^{p-1};\Z)$ jest $\Z$ dla $k=p$
i zerem poza tym stopniem, względna postać
Künnetha (bez członu $\operatorname{Tor}$, gdyż odpowiednia
grupa jest wolna) identyfikuje każdą summandę z
$H_q(F;\Z)$. Zmiana trywializacji działa na tej grupie
przez monodromię. Dla $q<0$ grupy są zerowe.

\emph{Krok 3: identyfikujemy $d^1$.}
Homomorfizm łączący pary $(E_p,E_{p-1})$ bierze granicę
względnego cyklu nad komórką i ogranicza ją do brzegu komórki.
Dla dwóch sąsiednich komórek suma współczynników
zawiera znak liczby incydencji ich orientacji, a
współczynnik z włókna trzeba przenieść wzdłuż drogi
łączącej wybrane punkty bazowe. Jest to dokładnie
różniczka komórkowych łańcuchów z lokalnym systemem
$\mathcal H_q(F)$, którego konstrukcję podaliśmy w
Sekcji~\ref{sec:lokalne-kolancuchy-v4}.
Stąd $E^2_{p,q}=H_p(B;\mathcal H_q(F))$.

\emph{Krok 4: różniczki wyższe i zbieżność.}
Z par kolejnych filtrów otrzymujemy dokładną parę,
a z niej wyższe różniczki i strony, jak w
Twierdzeniu~\ref{thm:pages-v5}.
Dla ustalonego stopnia całkowitego $n$ w pierwszej ćwiartce
możliwe są jedynie indeksy $0\le p\le n$,
zatem filtracja homologii ma skończoną długość.
Ponieważ $E_m=E$, filtracja łańcuchów jest skończona;
nie ma ukrytych klas pochodzących z nieskończonego
dopełnienia filtracji. Standardowa
identyfikacja cykli i brzegów stabilnych z dowodu
Twierdzenia~\ref{thm:convergence-v5} daje opis $E^\infty$.
Naturalność pochodzi z naturalności wszystkich etapów:
przeciwobrazów szkieletów, homomorfizmów łączących i
transportu wzdłuż dróg.
\end{proof}

\begin{uwaga}[Uwaga o zakresie dowodu]
Krok 2 korzysta z filtracyjnej wersji wycinania dla
CW-kompleksów, a nie ze zwykłego wycinania wzdłuż
domkniętych komórek bez sprawdzenia jego hipotez.
Dla wiązek lokalnie trywialnych odpowiedni izomorfizm
otrzymuje się najpierw po zastąpieniu $E_{p-1}$ małym
otwartym otoczeniem deformacyjnie retrahującym się
na $E_{p-1}$, a następnie przez wycinanie nad rozłącznymi
otwartymi $p$-komórkami. Uogólnienie na wszystkie fibracje
Serre'a wymaga odrębnego twierdzenia porównawczego.
\end{uwaga}

\subsection*{Od wiązek do ogólnych fibracji Serre'a}
Powyższy dowód przeprowadziliśmy dla lokalnie trywialnych wiązek.
Teraz uzasadnimy rozszerzenie na fibracje Serre'a bez zakładania
lokalnej trywialności. Rozróżnienie jest ważne: fibracja Serre'a
ma własność podnoszenia homotopii względem dysków, ale nie
musi posiadać lokalnych kart trywializujących.

\begin{lemat}[Porównanie z fibracją przestrzeni dróg]
\label{lem:serre-mappingpath-v6}
Niech $p:E\to B$ będzie surjektywną fibracją Serre'a, $B$ przestrzenią
łukowo spójną, a jej włókna niech mają typ homotopijny CW-kompleksów.
Zdefiniujmy
\[
E^h=\{(e,\gamma)\in E\times B^I:\gamma(0)=p(e)\},\qquad
p^h(e,\gamma)=\gamma(1).
\]
Przestrzeń dróg $B^I$ wyposażamy w topologię zwarto-otwartą,
a $I=[0,1]$. Wtedy $p^h$ jest fibracją Hurewicza, a mapa
$j:E\to E^h$, $j(e)=(e,c_{p(e)})$, jest mapą nad $B$.
Dla każdego $b\in B$ mapa włókien
$j_b:F_b\to (p^h)^{-1}(b)$ indukuje izomorfizmy
wszystkich grup homologii singularnej.
Co więcej, dla dowolnego podkompleksu CW $A\subseteq B$
mapa $j:p^{-1}(A)\to(p^h)^{-1}(A)$ indukuje
izomorfizmy homologii, a dla pary podkompleksów
$A'\subseteq A$ również homologii względnej.
\end{lemat}
\begin{proof}
Najpierw wyjaśnijmy własność Hurewicza. Dla homotopii
$h:Z\times I\to B$ oraz podniesienia początkowego
$\widetilde h_0(z)=(e_z,\gamma_z)$, gdzie
$\gamma_z(1)=h(z,0)$, pozostawiamy $e_z$ bez zmiany,
a drogę $\gamma_z$ przedłużamy od $h(z,0)$ do $h(z,t)$
fragmentem $s\mapsto h(z,st)$. Reparametryzując oba
odcinki liniowo, uzyskujemy ciągłe podniesienie $h$.

Mapa $E\to E^h$, $e\mapsto(e,c_{p(e)})$, jest
równoważnością homotopijną \emph{całych przestrzeni}: rzut
$r:E^h\to E$, $(e,\gamma)\mapsto e$, spełnia $rj=\id$,
a $jr\simeq\id$ przez skracanie drogi $\gamma$ od jej końca.
Ta homotopia nie zachowuje jednak rzutowania na $B$;
dlatego nie wystarcza do porównania poszczególnych włókien.

Dla $b\in B$ porównujemy długie ciągi homotopii
fibracji Serre'a $F_b\to E\to B$ i fibracji Hurewicza
$(p^h)^{-1}(b)\to E^h\to B$. W podstawie występuje
identyczność, a w przestrzeni całkowitej $j$ jest
równoważnością homotopijną. Z lematu pięciu w dodatnich
stopniach oraz z dokładności na poziomie zbiorów składowych
wynika, że $j_b$ jest słabą równoważnością homotopijną.
Słaba równoważność dowolnych przestrzeni indukuje
izomorfizm homologii singularnej: wynika to z twierdzenia
Whiteheada zastosowanego do ich kanonicznych przybliżeń CW
przez realizacje geometryczne kompleksów singularnych.

Teraz niech $A\subseteq B$ będzie podkompleksem.
Cofnięcia obu fibracji nad $A$ są fibracjami Serre'a;
porównanie ich długich ciągów homotopii na każdej składowej
$A$ daje słabą równoważność na przestrzeniach całkowitych
nad $A$. Stąd izomorfizm homologii. Jeżeli
$A'\subseteq A$, stosujemy do otrzymanych map par
długie ciągi homologii i lemat pięciu, co daje ostatnią tezę.
\end{proof}

\begin{lemat}[Lokalny rachunek nad komórką dla fibracji Hurewicza]
\label{lem:serre-cell-hurewicz-v6}
Niech $q:Y\to B$ będzie fibracją Hurewicza nad CW-kompleksem
$B$, a $Y_p=q^{-1}(B^p)$. Dla każdej zorientowanej
$p$-komórki $e^p_\alpha$ wybierzmy punkt $b_\alpha$ w jej
wnętrzu. Wtedy
\[
H_{p+q}(Y_p,Y_{p-1};\Z)\cong
\bigoplus_{\alpha\in\operatorname{Cells}_p(B)}
H_q(q^{-1}(b_\alpha);\Z),\qquad q\ge0,
\]
a grupy względne w stopniu $p+q$ z $q<0$ znikają.
Po tych identyfikacjach homomorfizm łączący
$H_{p+q}(Y_p,Y_{p-1})\to H_{p+q-1}(Y_{p-1},Y_{p-2})$
jest różniczką komórkową z transportem współczynników
$H_q(q^{-1}(b))$ wzdłuż dróg w podstawie.
\end{lemat}
\begin{proof}
Dla każdej mapy charakterystycznej
$\chi_\alpha:(D^p,S^{p-1})\to(B^p,B^{p-1})$
cofamy $q$ do fibracji nad dyskiem. Lemat
\ref{lem:serre-disk-v5} daje równoważność włóknistą
z produktem $D^p\times F_\alpha$; ograniczenie na
$S^{p-1}$ daje równoważność par. Dlatego homologia
względna cofnięcia wynosi
$H_{p+q}(D^p\times F_\alpha,S^{p-1}\times F_\alpha)
\cong H_q(F_\alpha)$, co wynika ze względnego wzoru
Künnetha, gdyż $H_p(D^p,S^{p-1})=\Z$.

Wyjaśnijmy przejście od pojedynczej komórki do całego
szkieletu. Wybierzmy w $B^p$ otoczenie $U$ szkieletu
$B^{p-1}$ deformacyjnie retrahujące się na $B^{p-1}$
względem tego szkieletu. Podnosząc tę deformację
z własności Hurewicza otrzymujemy równoważność
$q^{-1}(B^{p-1})\hookrightarrow q^{-1}(U)$.
Po wycięciu $q^{-1}(U)$ pozostają rozłączne części
pochodzące z wnętrz komórek. Zwarta nośność pojedynczego
łańcucha singularnego redukuje dowolny rachunek do
skończonej liczby komórek i daje sumę prostą.
Użyte wycinanie jest wycinaniem \emph{par po pogrubieniu}
skieletu, a nie nieuzasadnionym wycinaniem domkniętych
brzegów komórek.

Aby obliczyć $d^1$, wybieramy orientacje komórek i drogi
łączące punkty bazowe. Na ścianie dysku
charakterystycznego transportujemy klasę $H_q(F)$ do
włókna nad sąsiednią $(p-1)$-komórką. Wkład tej ściany
jest jej liczbą incydencji pomnożoną przez transport.
Dlaczego nie powstają inne współczynniki?
Na małym dysku w ścianie możemy przez homotopię mapy
charakterystycznej zastąpić jej zachowanie mapą o lokalnym
stopniu $+1$ lub $-1$; naturalność homomorfizmu
łączącego pary dyskowej mnoży wówczas generator
$H_p(D^p,S^{p-1})$ przez odpowiedni znak, pozostawiając
transportowaną klasę włókna. Sumowanie po ścianach
jest dokładnie definicją granicy komórkowej z lokalnymi
współczynnikami. Zależność od wyboru dróg zmienia
jedynie identyfikacje włókien systemu lokalnego,
nie sam kompleks.
\end{proof}

\begin{twierdzenie}[Ciąg spektralny Serre'a dla fibracji]
\label{thm:serre-general-v6}
Niech $p:E\to B$ będzie surjektywną fibracją Serre'a,
$B$ łukowo spójnym, skończeniewymiarowym CW-kompleksem,
a włókna niech mają typ homotopijny CW-kompleksów.
Dla wybranego $b_0\in B$ oznaczmy $F=p^{-1}(b_0)$.
Wtedy istnieje naturalny homologiczny ciąg spektralny
\[
E^2_{s,t}\cong H_s(B;\mathcal H_t(F;\Z)),\qquad
 d^r:E^r_{s,t}\longrightarrow E^r_{s-r,t+r-1},
\]
który silnie zbiega do $H_{s+t}(E;\Z)$, w tym sensie, że
$E^\infty_{s,t}\cong F_sH_{s+t}(E)/F_{s-1}H_{s+t}(E)$,
a filtracja jest skończona w każdym stopniu.
System $\mathcal H_t(F;\Z)$ wyznacza monodromia fibracji
na homologii włókien. Nie wymagamy lokalnej trywialności.
\end{twierdzenie}
\begin{proof}
Zastępujemy $p$ fibracją przestrzeni dróg $p^h:E^h\to B$
z Lematu~\ref{lem:serre-mappingpath-v6}. Rzutowanie $p^h$
ma własność podnoszenia Hurewicza, więc Lemat
\ref{lem:serre-cell-hurewicz-v6} oblicza stronę
$E^1_{s,t}=H_{s+t}((E^h)_s,(E^h)_{s-1})$
i identyfikuje $d^1$ z granicą komórkową z lokalnymi
współczynnikami $H_t((p^h)^{-1}(b))$.

Dla każdego szkieletu $B^s$ porównanie $j:E\to E^h$
indukuje izomorfizm na homologii jego przeciwobrazu;
na mocy długich ciągów par daje więc izomorfizm
$H_* (E_s,E_{s-1})\to H_*((E^h)_s,(E^h)_{s-1})$.
Są to izomorfizmy \emph{zgodne z homomorfizmami łączącymi},
ponieważ pochodzą z jednej mapy filtrowanych przestrzeni.
Dają zatem izomorfizm dokładnych par i wszystkich
stron ich ciągów spektralnych. Na włóknach $j$ indukuje
izomorfizmy homologii, więc identyfikuje systemy
lokalnych współczynników. Otrzymujemy
$E^2_{s,t}=H_s(B;\mathcal H_t(F;\Z))$.

Dla $d=\dim B<\infty$ filtracja
$E_{-1}=\varnothing\subseteq E_0\subseteq\cdots
\subseteq E_d=E$ jest ograniczona \emph{na poziomie
kompleksów łańcuchowych}. Stosujemy zatem
Twierdzenie~\ref{thm:convergence-v5}: daje ono stabilizację,
identyfikację kolejnych ilorazów i brak problemu
warunkowej zbieżności przy nieskończonej filtracji.
Naturalność zachodzi dla map fibracji respektujących
wybraną filtrację szkieletową; jej niezależność od wyboru
modelu CW wynika z porównania przez wspólne
przybliżenia komórkowe.
\end{proof}
\begin{uwaga}
Dla bazy o nieskończonym wymiarze także istnieje klasyczny
ciąg Serre'a w szerokiej ogólności. Przejście od
skończeniewymiarowych filtracji do sytuacji nieskończonej
wymaga osobnej analizy zbieżności; nie jest wnioskiem
z samego twierdzenia o filtracji ograniczonej.
\end{uwaga}


\section{Wersja kohomologiczna i transgresja}
\begin{twierdzenie}[Kohomologiczny ciąg spektralny]
Jeżeli $F\to E\to B$ jest lokalnie trywialną wiązką
nad spójnym skończeniewymiarowym CW-kompleksem $B$, to istnieje ciąg spektralny
kohomologiczny, dla którego
\[
E_2^{p,q}\cong H^p(B;\mathcal H^q(F)),\qquad
 d_r:E_r^{p,q}\longrightarrow E_r^{p+r,q-r+1},
\]
a $E_\infty^{p,q}$ są kolejnymi ilorazami odpowiedniej
filtracji $H^{p+q}(E;\Z)$. System $\mathcal H^q(F)$
ma włókna $H^q(F_b;\Z)$, zaś dla drogi $\gamma:b\to b'$
transport w tym systemie jest równy
$(T_{\bar\gamma})^*:H^q(F_b;\Z)\to H^q(F_{b'};\Z)$,
gdzie $\bar\gamma$ jest drogą odwrotną, a $T$ oznacza
transport włókien, rozpatrywany z dokładnością do homotopii.
Odwrócenie drogi wynika z kontrawariantności kohomologii.
W przypadku trywialnej monodromii
system $\mathcal H^q(F)$ jest stały. Dla orientowalnej
wiązki sferowej najważniejszą różniczką jest transgresja.
\end{twierdzenie}
\begin{proof}
Niech $B^p$ będzie $p$-szkieletem, $E_p$ jego przeciwobrazem
w $E$, a $m=\dim B$. W kompleksie kołańcuchów singularnych
wprowadzamy malejącą filtrację
\[
 F^pC^n(E;\Z)=
 \ker\bigl(C^n(E;\Z)\longrightarrow C^n(E_{p-1};\Z)\bigr).
\]
Dla $p=0$ mamy $F^0C^*=C^*(E)$, a dla $p=m+1$
otrzymujemy $F^{m+1}C^*=0$; koróżniczka zachowuje
filtrację. Ograniczenie kołańcucha z $E$ do $E_p$ daje
izomorfizm kompleksów
\[
 F^pC^*(E)/F^{p+1}C^*(E)
 \cong C^*(E_p,E_{p-1};\Z).
\]
Rzeczywiście, każdy kołańcuch na $E_p$, znikający na
$E_{p-1}$, przedłuża się do $E$ przez przypisanie zera
pozostałym sympleksom singularnym; nie nakładamy warunku
skończonego nośnika na kołańcuchy.

Strona pierwsza jest więc równa
$E_1^{p,q}=H^{p+q}(E_p,E_{p-1};\Z)$.
Obliczenie względne nad komórkami podstawy identyfikuje
te grupy z komórkowymi kołańcuchami o współczynnikach
w systemie $\mathcal H^q(F)$: dla nieskończenie wielu
komórek w ustalonym wymiarze występuje \emph{produkt},
nie suma bezpośrednia. Homomorfizm łączący pary daje
$d_1$, czyli koróżniczkę komórkową z monodromią.
Otrzymujemy $E_2^{p,q}=H^p(B;\mathcal H^q(F))$.
Wyższe różniczki mają bistopień $(r,1-r)$ na mocy
konstrukcji kolejnych stron dokładnej pary.
Ponieważ filtracja jest ograniczona, a argument
stabilizacji cykli i brzegów jest dualny do dowodu
Twierdzenia~\ref{thm:convergence-v5}, dostajemy
deklarowaną filtrację kohomologii $E$ i jej ilorazy.
\end{proof}
\begin{definicja}[Klasa transgresywna]
Niech $q\ge1$ i niech $u\in E_2^{0,q}$ będzie
klasą w kohomologicznym ciągu spektralnym.
Nazywamy ją \emph{transgresywną}, jeśli
przeżywa do strony $E_{q+1}^{0,q}$, to znaczy
kolejne wychodzące z niej różniczki
$d_2,\ldots,d_q$ są zerowe. Wówczas
$d_{q+1}u\in E_{q+1}^{q+1,0}$ jest jej transgresją
(przy czym przez $u$ oznaczamy też klasę reprezentowaną
przez nią na stronie $q+1$).
Nie każdy element $H^q(F)$ jest transgresywny:
wcześniejsza różniczka może go unicestwić.
\end{definicja}
\begin{uwaga}[Iloczyn i reguła Leibniza]
Jeśli $B$ jest jednospójna, włókno ma odpowiedni typ CW,
a współczynniki należą do przemiennego pierścienia,
ciąg kohomologiczny ma strukturę multiplikatywną;
dla jednorodnego $a\in E_r^{p,q}$ i $b$ zachodzi
$d_r(ab)=d_r(a)b+(-1)^{p+q}a\,d_r(b)$.
Wymaga to filtracyjnego odwzorowania diagonalnego
i konstrukcji iloczynu na stronie zerowej; jest ono zgodne
z iloczynem kubkowym całkowitej kohomologii. Nie wolno
zakładać identycznego wzoru przy lokalnym systemie o
nieprzemiennym działaniu bez zdefiniowania odpowiedniego
sparowania systemów współczynników.
\end{uwaga}

\begin{wniosek}[Kohomologiczna wersja dla fibracji Serre'a]
Przy hipotezach Twierdzenia~\ref{thm:serre-general-v6}
ciąg spektralny z poprzedniego twierdzenia istnieje także
dla fibracji Serre'a. Jego strona druga ma postać
$E_2^{p,q}=H^p(B;\mathcal H^q(F;\Z))$ i zbiega
do $H^{p+q}(E;\Z)$ przez skończoną filtrację.
\end{wniosek}
\begin{proof}
Porównanie z fibracją przestrzeni dróg z
Lematu~\ref{lem:serre-mappingpath-v6} daje izomorfizm
kohomologii wszystkich par przeciwobrazów szkieletów.
Po odwróceniu strzałek w ich długich ciągach dokładnych
otrzymujemy izomorfizm kohomologicznych dokładnych par.
Obliczenie komórkowe jest dualne do
Lematu~\ref{lem:serre-cell-hurewicz-v6}; przy nieskończonym
zbiorze komórek w danym wymiarze sumę łańcuchów
zastępuje iloczyn kołańcuchów, jak wymaga definicja
kohomologii komórkowej. Skończeniewymiarowość bazy
zapewnia skończoną filtrację i zbieżność.
\end{proof}

\section{Przykład podstawowy: fibracja Hopfa i niezerowa \texorpdfstring{$d_2$}{d2}}
Rozważmy wiązkę $S^1\to S^3\to S^2$ z
Twierdzenia~\ref{thm:hopf-v4}. Podstawa jest jednospójna,
więc monodromia jest trywialna. Na stronie drugiej
mamy tylko cztery niezerowe pozycje:
$E^2_{0,0}=E^2_{2,0}=E^2_{0,1}=E^2_{2,1}=\Z$.

\begin{center}
\begin{tikzpicture}[x=1.25cm,y=1.0cm,>=Stealth]
\draw[->] (-.35,0)--(3.3,0) node[right]{$p$};
\draw[->] (0,-.25)--(0,2.3) node[above]{$q$};
\foreach \i in {0,1,2,3}{\draw(\i,.07)--(\i,-.07) node[below] {\small\i};}
\foreach \i in {0,1,2}{\draw(.07,\i)--(-.07,\i) node[left] {\small\i};}
\foreach \x/\y in {0/0,2/0,0/1,2/1}{\node[draw,rounded corners,fill=blue!7,minimum width=.65cm,minimum height=.43cm] at (\x,\y){$\Z$};}
\draw[->,very thick,red!70!black] (1.75,.13)--(.25,.88) node[midway,above] {$d^2$};
\node[align=left] at (2.45,1.64){\footnotesize $E^2_{p,q}$\\[-1pt]\footnotesize fibracji Hopfa};
\end{tikzpicture}
\end{center}

\begin{twierdzenie}
W homologicznym ciągu spektralnym Hopfa homomorfizm
$d^2:E^2_{2,0}\to E^2_{0,1}$ jest izomorfizmem
mnożenia przez $\pm1$. W konsekwencji na stronie
$E^\infty$ pozostają tylko grupy w pozycjach $(0,0)$
i $(2,1)$, obie izomorficzne z $\Z$.
\end{twierdzenie}
\begin{proof}
Na stronie $E^2$ nie ma niezerowych różniczek poza
wskazaną: różniczka ma bistopień $(-r,r-1)$,
a podstawa ma tylko stopnie $0$ i $2$.
Pozycja $(2,0)$ odpowiada stopniowi całkowitemu $2$,
a $(0,1)$ stopniowi $1$. Ponieważ
$H_1(S^3;\Z)=H_2(S^3;\Z)=0$, obie pozycje muszą
zniknąć na stronie $E^\infty$; jedyną możliwością
jest $d^2$ o zerowym jądrze i zerowym kokernelu.
Jako endomorfizm $\Z$ jest więc mnożeniem przez $\pm1$.
Pozostałe dwa wyrazy nie mogą zostać zabite ani nie mogą
być obrazami różniczki z pierwszej ćwiartki.
Na stopniach $0$ i $3$ nie ma problemu rozszerzeń,
co odtwarza homologię $S^3$.
\end{proof}
\begin{uwaga}
W kohomologii ta sama fibracja daje
$d_2:E_2^{0,1}\to E_2^{2,0}$, również izomorfizm.
Klasa generatora $H^1(S^1;\Z)$ transgresuje do generatora
$H^2(S^2;\Z)$. Nie należy zamieniać kierunków strzałek
homologicznej i kohomologicznej różniczki.
\end{uwaga}

\section{Przykład z nietrywialną monodromią: butelka Kleina}
Niech $r:S^1\to S^1$ będzie odbiciem, a $K$ jego
\emph{torusem odwzorowania}, czyli przestrzenią ilorazową
$(S^1\times[0,1])/{\sim}$, gdzie $(z,1)\sim(r(z),0)$.
Jest to butelka Kleina opisana w rozdziale o fibracjach.
Rzut na współrzędną przedziału modulo sklejenie daje
wiązkę $S^1\to K\to S^1$. Działanie monodromii
na $H_0(S^1;\Z)=\Z$ jest trywialne, a na
$H_1(S^1;\Z)=\Z$ jest mnożeniem przez $-1$.
Dla współczynników $\Z_-$ otrzymujemy komórkowy kompleks
$\Z\xrightarrow{-2}\Z$, a dla stałych $\Z$
kompleks $\Z\xrightarrow{0}\Z$. Zatem
\[
E^2_{0,0}=\Z,\quad E^2_{1,0}=\Z,\quad
E^2_{0,1}=\Z/2,\quad E^2_{1,1}=0.
\]
Ponieważ $\dim S^1=1$, wszystkie $d^r$ dla $r\ge2$
mają zerowe cele. Dostajemy $H_0(K;\Z)=\Z$,
$H_2(K;\Z)=0$ i krótki ciąg dokładny
$0\to\Z/2\to H_1(K;\Z)\to\Z\to0$.
Rozszczepia się ona, bo $\Z$ jest wolna, zatem
$H_1(K;\Z)\cong\Z\oplus\Z/2$.
Wyrzucenie systemów lokalnych ze wzoru $E^2$ dałoby
błędny wynik $H_1(K)\cong\Z^2$.

\section{Wiązki sferowe, klasa Eulera i ciąg Gysina}
\begin{definicja}[Wiązka sferowa zorientowanej wiązki]
Niech $\xi\to B$ będzie orientowalną rzeczywistą wiązką
wektorową rzędu $m+1$, wyposażoną w metrykę włóknową.
Niech $D(\xi)$ i $S(\xi)$ oznaczają wiązki odpowiednio
jednostkowych dysków i jednostkowych sfer. Mamy
wiązki $S^m\to S(\xi)\xrightarrow{\pi}B$ i
$D^{m+1}\to D(\xi)\to B$.
\end{definicja}
\begin{twierdzenie}[Izomorfizm Thoma i klasa Eulera]
\label{thm:thom-v5}
Niech $B$ będzie skończonym CW-kompleksem, a
$\xi\to B$ zorientowaną rzeczywistą wiązką wektorową
rzędu $r\ge1$ z metryką włóknową. Oznaczmy przez
$D(\xi)$ i $S(\xi)$ podwiązki jednostkowych dysków
oraz sfer, a przez $\pi:D(\xi)\to B$ rzutowanie.
Wówczas istnieje dokładnie jedna klasa
$u_\xi\in H^r(D(\xi),S(\xi);\Z)$, której ograniczenie
do każdej pary włókien $(D^r_b,S^{r-1}_b)$ jest
jej dodatnim generatorem wyznaczonym przez orientację
$\xi_b$. Dla każdego $k\ge0$ homomorfizm
\[
T_\xi:H^k(B;\Z)\longrightarrow
H^{k+r}(D(\xi),S(\xi);\Z),\qquad
\alpha\longmapsto\pi^*\alpha\smile u_\xi
\]
jest izomorfizmem. Obraz $u_\xi$ przez
$H^r(D(\xi),S(\xi))\to H^r(D(\xi))\cong H^r(B)$
nazywamy klasą Eulera $e(\xi)$.
\end{twierdzenie}
\begin{proof}
\emph{Krok 1: filtracja pary wiązek.}
Dla szkieletu $B^p$ oznaczmy
$D_p=D(\xi)|_{B^p}$, $S_p=S(\xi)|_{B^p}$.
Relatywny kompleks kołańcuchów pary $(D,S)$
filtrujemy przez jądra ograniczeń do
$(D_{p-1},S_{p-1})$. Różniczki kolejnych stron
tej filtracji pochodzą z długich ciągów dokładnych
trójek $(D_p,D_{p-1}\cup S_p,S_p)$; całkowita
kohomologia jest $H^*(D,S)$, ponieważ $B$ ma
skończony wymiar.

\emph{Krok 2: wyrazy względne nad komórkami.}
Dla zorientowanej $p$-komórki $e^p_\alpha$ cofnięcie
wiązki na dysk charakterystyczny jest trywialne.
Wycinanie po pogrubieniu brzegu dysku i względny
iloczyn Künnetha dają
\[
\begin{aligned}
&H^{p+q}(D_p,D_{p-1}\cup S_p;\Z)\\
&\qquad\cong
\bigoplus_{\alpha\in\operatorname{Cells}_p(B)}
H^p(D^p,S^{p-1};\Z)\otimes
H^q(D^r,S^{r-1};\Z).
\end{aligned}
\]
Wszystkie grupy $H^q(D^r,S^{r-1};\Z)$ znikają
poza $q=r$, gdzie są równe $\Z$. Orientacja $\xi$
sprawia, że transport tego generatora wokół każdej
pętli jest tożsamością. Otrzymany \emph{jedyny}
niezerowy wiersz strony pierwszej jest więc
komórkowym kompleksem kołańcuchów
$C^p_{\mathrm{CW}}(B;\Z)$, przesuniętym w stopniu o $r$.
Wobec tego
\[
E_2^{p,q}=\begin{cases}H^p(B;\Z),&q=r,\\0,&q\ne r.
\end{cases}
\]
Żadna różniczka $d_s$ dla $s\ge2$ nie może łączyć
niezerowych wierszy. Ciąg spektralny stabilizuje się
na $E_2$. Dla każdego całkowitego stopnia $k+r$
pozostaje dokładnie jeden możliwy iloraz filtracji,
zatem nie powstaje także problem rozszerzeń.
W rezultacie otrzymujemy naturalne \emph{addytywne}
izomorfizmy
$\Theta_k:H^k(B;\Z)\xrightarrow{\cong}H^{k+r}(D,S;\Z)$.

\emph{Krok 3: konstrukcja globalnej klasy.}
Niech $1_B\in H^0(B;\Z)$ będzie klasą funkcji
przyjmującej wartość $1$ na każdej składowej bazy.
Definiujemy $u_\xi=\Theta_0(1_B)$.
Po ograniczeniu filtracji do dowolnej komórki
powraca wybrany dodatni generator względnej kohomologii
włókna. Ponieważ ograniczenia do włókien są zgodne
z transportem indukowanym przez orientację,
klasa $u_\xi$ ma żądaną własność. Jeżeli druga klasa
ma te same ograniczenia, jej różnica odpowiada
poprzez $\Theta_0$ klasie z $H^0(B)$ równej zeru
na każdej składowej, jest więc zerowa.

\emph{Krok 4: identyfikacja izomorfizmu z iloczynem.}
Zdefiniujmy $T_\xi(\alpha)=\pi^*\alpha\smile u_\xi$.
Filtracja komórkowa jest zgodna z iloczynem kubkowym:
na ilorazie odpowiadającym komórce $e^p$ mapa
$T_\xi$ jest zewnętrznym iloczynem
\[
H^p(D^p,S^{p-1};\Z)\longrightarrow
H^{p+r}(D^p\times D^r,
 (S^{p-1}\times D^r)\cup(D^p\times S^{r-1});\Z)
\]
z dodatnim generatorem $H^r(D^r,S^{r-1};\Z)$.
Obie grupy są izomorficzne z $\Z$, a mapa wysyła
generator na generator, jest więc izomorfizmem.
Przez naturalność iloczynu otrzymujemy izomorfizm
na wszystkich składnikach strony pierwszej oraz
przemienność z różniczką $d_1$.
Porównanie filtracji skończonych (indukcyjnie z lematu
pięciu dla kolejnych krótkich ciągów filtracji)
pokazuje, że $T_\xi$ jest izomorfizmem w każdym stopniu.
W szczególności $T_\xi(1_B)=u_\xi$, co ustala zgodność
w stopniu zero. Dla pozostałych stopni nie wystarczy
sama ta równość: izomorfizm $T_\xi$ został dowiedziony
osobno przez porównanie jego działania na kolejnych
ilorazach filtracji. Nie potrzebujemy twierdzić, że
wcześniej otrzymane addytywne identyfikacje $\Theta_k$
są z nim automatycznie identyczne.
\end{proof}
\begin{uwaga}[Wiązka nieorientowalna i współczynniki lokalne]
Jeżeli $\xi$ nie jest orientowalna, drugi krok nie daje
stałego wiersza $\Z$: otrzymujemy lokalny system orientacji
$\mathcal O_\xi$ i odpowiednią wersję izomorfizmu Thoma
z lokalnymi współczynnikami. Dlatego w zwykłym
całkowitoliczbowym twierdzeniu powyżej orientowalność
nie jest założeniem zbędnym.
\end{uwaga}


\begin{twierdzenie}[Ciąg Gysina]
\label{thm:gysin-v5}
Dla zorientowanej wiązki sferowej nad skończonym CW-kompleksem
jak powyżej istnieje
długi ciąg dokładny
\[
\begin{aligned}
\cdots &\to H^{k-m-1}(B;\Z)
\xrightarrow{\smile e(\xi)}H^k(B;\Z)
\xrightarrow{\pi^*}H^k(S(\xi);\Z)\\
&\xrightarrow{\pi_!}H^{k-m}(B;\Z)
\xrightarrow{\smile e(\xi)}H^{k+1}(B;\Z)\to\cdots.
\end{aligned}
\]
Jeżeli $m\ge1$, to w kohomologicznym ciągu spektralnym
Serre'a generator $H^m(S^m;\Z)$, jeśli jest transgresywny,
przechodzi pod $d_{m+1}$ w klasę $e(\xi)$
(z dokładnością do znaku). Dla $m=0$ włókno $S^0$
nie jest spójne i powyższy opis transgresji nie ma
tej postaci; sam ciąg Gysina nadal wynika z pary
$(D(\xi),S(\xi))$.
\end{twierdzenie}
\begin{proof}
Stosujemy długi ciąg dokładny pary
$(D(\xi),S(\xi))$:
$H^k(D,S)\to H^k(D)\to H^k(S)\to H^{k+1}(D,S)$.
Ponieważ $D(\xi)$ deformacyjnie retrahuje się na przekrój
zerowy, mamy $H^k(D)\cong H^k(B)$.
Z Twierdzenia~\ref{thm:thom-v5} mamy
$H^k(D,S)\cong H^{k-m-1}(B)$ oraz
$H^{k+1}(D,S)\cong H^{k-m}(B)$.
Mapa $H^k(D,S)\to H^k(D)$ jest przez naturalność
iloczynu kubkowego mnożeniem przez obraz klasy Thoma,
czyli przez $e(\xi)$. Homomorfizm $\pi_!$ jest
połączeniem mapy łączącej i odwrotności izomorfizmu Thoma.
To daje wszystkie strzałki oraz dokładność.

Dla $m\ge1$ w ciągu spektralnym wiązki sferowej możliwe
są jedynie wiersze $q=0,m$. Na pierwszej stronie, na której
różniczka może łączyć te wiersze, jej stopień wynosi $m+1$.
Porównanie homomorfizmu łączącego pary dyskowej
z homomorfizmem z dokładnej pary filtracji identyfikuje
transgresję generatora włókna z obrazem klasy Thoma.
Zgodność znaku zależy od przyjętych orientacji komórek.
\end{proof}
\begin{przyklad}
Dla trywialnej wiązki $\xi$ istnieje nigdzie niezerowy
przekrój, więc $e(\xi)=0$, a ciąg Gysina rozszczepia
poszczególne grupy (niekoniecznie w sposób kanoniczny).
Dla wiązki Hopfa $S^1\to S^3\to S^2$ klasa Eulera
odpowiedniej zorientowanej wiązki rzeczywistej rzędu $2$
ma wartość $\pm1$ razy generator $H^2(S^2;\Z)$,
co jest równoważne izomorfizmowi $d_2$ w poprzednim rachunku.
\end{przyklad}

\section{Zastosowanie do przestrzeni pętli i ograniczenia metody}
Dla $n\ge2$ mamy fibrację przestrzeni dróg
$\Omega S^n\to PS^n\to S^n$, której przestrzeń całkowita
jest ściągalna. W tym przykładzie korzystamy z \emph{wersji dla
fibracji Hurewicza} z
Twierdzenia~\ref{thm:serre-general-v6}:
fibracja przestrzeni dróg nie została tutaj wykazana
jako lokalnie trywialna wiązka.
Jej podstawa $S^n$ jest skończeniewymiarowym
CW-kompleksem, a włókno $\Omega S^n$
ma typ homotopijny CW-kompleksu, co wynika
np. z przybliżenia CW dla przestrzeni pętli
(zgodnego z grupami homotopii).
Dlatego wszystkie hipotezy przywołanego
twierdzenia są spełnione. Poniższy rachunek
jest zastosowaniem wersji \emph{już udowodnionej},
nie dodatkowym dowodem jej ogólności. Ponieważ $S^n$ jest jednospójna,
lokalne współczynniki są trywialne. Na stronie drugiej
występują tylko kolumny $p=0,n$.
Różniczka mogąca unicestwić klasę
$E^2_{n,0}=H_n(S^n;\Z)\cong\Z$ ma stopień $r=n$
i cel $E^n_{0,n-1}$. Wszystkie wcześniejsze różniczki
z tej pozycji są zerowe z powodów wymiarowych,
a pozycja $(n,0)$ nie może być trafiona przez różniczkę
z pierwszej ćwiartki o $p>n$. Skoro $H_n(PS^n)=0$,
$d^n:E^n_{n,0}\to E^n_{0,n-1}$ musi być iniekcją.
Również $H_{n-1}(PS^n)=0$, a $(0,n-1)$ może być
trafiona tylko w ten sposób, więc ta mapa jest surjekcją.
Zatem $H_{n-1}(\Omega S^n;\Z)\cong\Z$.
Jest to zgodne z izomorfizmem
$\pi_{n-1}(\Omega S^n)\cong\pi_n(S^n)$, ale wynika tu
z rachunku homologicznego.

\begin{uwaga}[Dokładny zakres otrzymanych wyników]
Niniejszy rozdział konstruuje wszystkie strony i różniczki
filtrowanego kompleksu oraz dowodzi ciągu Serre'a
dla wiązek lokalnie trywialnych
(Twierdzenie~\ref{thm:serre-v5}).
\emph{Ponadto} przez porównanie z fibracją przestrzeni
dróg i lokalne obliczenia dla fibracji Hurewicza
udowodniliśmy wersję dla fibracji Serre'a
o włóknach typu CW i \emph{skończeniewymiarowej}
łukowo spójnej bazie CW
(Twierdzenie~\ref{thm:serre-general-v6}).
Nie należy więc opisywać całej wersji dla fibracji
Serre'a jako nieudowodnionej.

Osobnej analizy wymagają natomiast:
ogólne twierdzenia o silnej zbieżności przy
\emph{nieograniczonej} filtracji, pełna konstrukcja
multiplikatywności z dowolnymi skręconymi
współczynnikami i formuły transgresji poza
przeanalizowanymi przypadkami wiązek sferowych.
Wyniki te nie są automatycznymi konsekwencjami
samego zapisu $E_2^{p,q}=H^p(B;H^q(F))$.
\end{uwaga}

\section{Podsumowanie metody obliczeniowej}
Przy konkretnej fibracji ustalamy najpierw system lokalny
$\mathcal H_q(F)$ i rzeczywiste działanie monodromii.
Następnie obliczamy $H_p(B;\mathcal H_q(F))$ oraz
rozrysowujemy niezerowe pozycje strony $E^2$.
Sprawdzamy, które różniczki mogą być niezerowe z uwagi
na ich bistopień i znane grupy przestrzeni całkowitej.
Dopiero po obliczeniu $E^\infty$ rozwiązujemy problem
rozszerzeń. Wersja kohomologiczna dostarcza dodatkowo
iloczynów oraz transgresji, której podstawowym przykładem
jest klasa Eulera zorientowanej wiązki sferowej.
\clearpage


